% ADDED figure (2026-11-23): the Klein-Nishina distribution in the lab frame for omega/m -> 0, 0.2, 1, 5 (after the figure of the user's notes); generated by qft-C/tools/l22/figs/make_figs.sh
\tikzset{
  fdl/.style={line width=0.8pt, black},
  fdx/.style={circle, fill=black, inner sep=0pt, minimum size=3.6pt},
  fdv/.style={circle, fill=figB, draw=figB, inner sep=0pt, minimum size=5pt},
  fdlab/.style={font=\scriptsize, inner sep=2.6pt},
  fdmom/.style={->, >=Stealth, figO, line width=0.6pt},
}
\tikzset{
  fda/.style={fdl, postaction={decorate}, decoration={markings, mark=at position 0.55 with {\arrow{Stealth[length=5pt]}}}},
  fdS/.style={font=\small},
}
\tikzset{
  fline/.style={line width=0.8pt, black, postaction={decorate}, decoration={markings, mark=at position 0.56 with {\arrow{Stealth[length=5.5pt, width=4.5pt]}}}},
  sline/.style={line width=0.8pt, black, dash pattern=on 3.5pt off 2.2pt},
  blob/.style={draw, line width=0.8pt, fill=figB!12},
}
\tikzset{
  pline/.style={line width=0.8pt, black, decorate, decoration={snake, amplitude=1.6pt, segment length=5.2pt}},
}
\begin{tikzpicture}
\begin{axis}[nfig, width=9.6cm, height=6.6cm, xmin=0, xmax=180, ymin=0, ymax=2.15, xlabel={scattering angle $\theta$ (degrees), lab frame}, ylabel={$\dfrac{m^2}{\pi\alpha^2}\,\dfrac{d\sigma}{d\cos\theta}$}, xtick={0,30,60,90,120,150,180}, ytick={0,0.5,1,1.5,2}, clip=false]
  \addplot[black, domain=0:180, samples=160] {1+cos(x)^2};
  \addplot[figB, domain=0:180, samples=160] {(1/(1+0.2*(1-cos(x))))^2*(1/(1+0.2*(1-cos(x)))+1+0.2*(1-cos(x))-sin(x)^2)};
  \addplot[figO, domain=0:180, samples=160] {(1/(1+1*(1-cos(x))))^2*(1/(1+1*(1-cos(x)))+1+1*(1-cos(x))-sin(x)^2)};
  \addplot[figR, domain=0:180, samples=160] {(1/(1+5*(1-cos(x))))^2*(1/(1+5*(1-cos(x)))+1+5*(1-cos(x))-sin(x)^2)};
  \node[font=\scriptsize, anchor=west] at (axis cs:182,2.0) {$\omega/m \to 0$: Thomson, $1 + \cos^2\theta$};
  \node[font=\scriptsize, anchor=west, figB] at (axis cs:182,1.08) {$\omega/m = 0.2$};
  \node[font=\scriptsize, anchor=west, figO] at (axis cs:182,0.40) {$\omega/m = 1$};
  \node[font=\scriptsize, anchor=west, figR] at (axis cs:182,0.09) {$\omega/m = 5$};
  \node[font=\scriptsize, gray, anchor=south west] at (axis cs:3,2.02) {forward value $2$ at every energy};
\end{axis}
\end{tikzpicture}
